\documentclass[11pt]{article}
\usepackage{preamble}
\renewcommand{\answer}[1]{}

\begin{document}

\ladderheading{AP Statistics \textperiodcentered\ Unit 4 \textperiodcentered\ Topic 4.9 \textperiodcentered\ B: 2027-01-28 \textperiodcentered\ E: 2027-03-17}{Two Test Plans for Cure Time}

\focusbanner{TODAY'S PROBLEM}

Suppose line N uses a new formula for sealant, a material that closes gaps. Line S uses the standard formula. A team takes independent simple random samples of 18 of 600 recent N batches and 18 of 540 recent S batches. Mean cure times are 41.2 and 44.0 minutes; the boxplots are shown. Before sampling, the team asked, ``Do the two populations differ in mean cure time?'' Formula was not randomly assigned to lines.\par\smallskip \textbf{Rhea:} ``Keep the original question and use population means. A difference between lines may have other causes.''\par \textbf{Sol:} ``Since $41.2<44.0$, use a lower tail with sample means. A significant result proves the formula caused faster curing.''\par
\begin{center}
\begin{tikzpicture}
\begin{axis}[
  width=0.61\linewidth,
  height=1.44in,
  ytick={2,1}, yticklabels={{Line N},{Line S}}, ymin=0.4, ymax=2.6,
  axis x line=bottom, axis y line*=left, y axis line style={draw=none}, boxplot/box extend=0.5, xmajorgrids, enlarge x limits=0.06,
  xlabel={Cure time (min)}, tick label style={font=\small}, label style={font=\small},
  ytick style={draw=none}, yticklabel style={font=\small\bfseries}
]
\addplot[boxplot prepared={lower whisker=34, lower quartile=38.5, median=41, upper quartile=44, upper whisker=48.5, draw position=2}, fill=calloutblue, draw=dayblue] coordinates {};
\addplot[boxplot prepared={lower whisker=35.5, lower quartile=40, median=44, upper quartile=48, upper whisker=52.5, draw position=1}, fill=calloutblue, draw=dayblue] coordinates {};
\end{axis}
\end{tikzpicture}
\end{center}
\begin{firsttakebox}{FIRST TAKE --- before the video (5 min)}
Did the team originally ask whether N was faster or whether the lines differed in either direction? Explain in one sentence.
\workspace{0.9in}
\end{firsttakebox}

\begin{callout}{\IconBook\ HYPOTHESES FOR TWO MEANS (keep on the page)}{calloutgreen}
For two independent groups with a quantitative response, use a two-sample $t$ test for
$\mu_1-\mu_2$. The null states no population difference: $H_0:\mu_1-\mu_2=0$.
Choose $<$ or $>$ only when the research question specified that direction before seeing
the data; otherwise use $\neq$. Hypotheses use population parameters, never observed sample
means. Check independent random samples or random assignment, both 10 percent conditions
when sampling without replacement, and both sample shapes when either $n<30$. Random sampling
supports population scope; random assignment is needed to isolate a causal treatment effect.
\end{callout}

\newpage
\textbf{Finish after the video:}\par\smallskip

\dokbadge{1}~\textbf{(a)} Define $\mu_N$ and $\mu_S$ in context.\par
\workspace{0.7in}

\dokbadge{2}~\textbf{(b)} State $H_0$ and $H_a$ in N-minus-S order to match the question asked before sampling.\par
\workspace{1.4in}

\dokbadge{3}~\textbf{(c)} In 3--4 sentences, judge the two plans. Explain why the original question sets the tail, why hypotheses use population means, and why these two lines cannot isolate the formula as the cause.\par
\workspace{2.0in}

\begin{sentenceframebox}\raggedright\textbf{Frame:} The original question asks \blank[1.8in], so the alternative should \blank[1.8in].\end{sentenceframebox}

\begin{sentenceframebox}\raggedright\textbf{Frame:} The hypotheses describe \blank[1.8in]. A causal claim is limited because \blank[2.0in].\end{sentenceframebox}

\begin{summaryexitbox}{EXIT --- turn this in}
One thing the video changed about my first take: \hrulefill\par
\workspace{0.4in}
\end{summaryexitbox}

\end{document}
